\scanpage{285}
\mspage{235}

dans $\unsure{\mathscr X|\mathcal X}_Q$, qui intègre le diagramme
\diagram{\begin{tikzcd}[ampersand replacement=\&,row sep=tiny]
 \& X_1\\
 X_0 \arrow[ur] \arrow[dr] \&\\
 \& X_2.
\end{tikzcd}}
Pour que $Q$ soit $W$-cocartésien, il f. et s. que l'inclusion
\[
i:\mathscr X\longrightarrow\mathscr X_Q
\]
soit \deleted{dans} $W$.

Comme $\mathscr X_Q$ est lisse sur $\widetilde Q$, et 3 est objet initial de $\widetilde Q$, il s'ensuit que l'inclusion de la fibre de \inserted{$\mathscr X_Q$ en le sommet} $3$ \deleted{de} $\unsure{\widetilde Q|Q}$ est dans $W$
\[
 Y\xrightarrow{i_Y}\deleted{\unsure{Q}}\mathscr X_Q\qquad\in W.
\]
Considérons \unsure{donc|alors} le diagramme
\diagram{\begin{tikzcd}[ampersand replacement=\&]
 \mathscr X \arrow[rr,"{\varphi_Q}"] \arrow[dr,"i"',""{name=left,coordinate,pos=.53}] \& \& Y \arrow[dl,"{i_Y\in W}", ""{name=right,coordinate,pos=.53}]\\
 \& \mathscr X_Q \arrow[from=right,to=left,dashed,"{\alpha}"'] \&
\end{tikzcd}}
Ce diagramme n'est pas commutatif, mais il y a une flèche de commutativité
\[
 \alpha:i_Y\varphi_Q\longrightarrow i.
\]
On sait que ceci implique que \deleted{$i\in W$}
\[
 i\in W \Longleftrightarrow (i_Y\varphi_Q)\in W
\]
(cf lemme p. 2), et comme $i_Y\in W$, \deleted{$i_Y\varphi_Q\in W\Longleftrightarrow$}
\[
 i_Y\varphi_Q\in W\Longleftrightarrow\varphi_Q\in W,
\]
d'où
\[
 i\in W\Longleftrightarrow\varphi_Q\in W\qquad \text{qed.}
\]
Nous pouvons maintenant prouver

% archive sheet number: 285

